\section{2026年高考题}
\subsection{2026年全国I卷}\label{gk:2026年全国I卷}

\begin{description}
    \item[一、选择题] 本题共8小题，每小题5分，共40分. 在每小题给出的四个选项中，只有一项是符合题目要求的. 
    \begin{enumerate}[leftmargin=5pt]
        \item 样本数据 $6,8,4,5,12$ 的中位数为

        {\centering\begin{tabular*}{\linewidth}{@{}@{\extracolsep{\fill}}llll@{}}
          \(\text{A. }5\)  & \(\text{B. }6\)  & \(\text{C. }8\)  & \(\text{D. }9\)
        \end{tabular*}\par}

        \item 已知平面向量 $\bm a,\bm b$ 不共线，且 $2\bm a+y\bm b=x\bm a-3\bm b$，则

        {\centering\begin{tabular*}{\linewidth}{@{}@{\extracolsep{\fill}}llll@{}}
          \(\text{A. }x=2,y=-3\)  & \(\text{B. }x=-2,y=3\)  & \(\text{C. }x=2,y=3\)  & \(\text{D. }x=-2,y=-3\)
        \end{tabular*}\par}

        \item 已知集合 $A=\{\sin\frac{7\pi}{6},\cos\frac{5\pi}{3},\tan\frac{5\pi}{4}\}$，$B=\{-\frac{3}{2},-\frac{1}{2},1\}$，则 $A\cap B=$

        {\centering\begin{tabular*}{\linewidth}{@{}@{\extracolsep{\fill}}llll@{}}
          \(\text{A. }\{-\frac{3}{2},-\frac{1}{2}\}\)  & \(\text{B. }\{-\frac{3}{2},1\}\)  & \(\text{C. }\{-\frac{1}{2},1\}\)  & \(\text{D. }\{-\frac{3}{2},-\frac{1}{2},1\}\)
        \end{tabular*}\par}

        \item 曲线 $y=5x+8\ln x$ 在点 $(1,5)$ 处的切线方程为

        {\centering\begin{tabular*}{\linewidth}{@{}@{\extracolsep{\fill}}llll@{}}
          \(\text{A. }y=3x+2\)  & \(\text{B. }y=5x\)  & \(\text{C. }y=8x-3\)  & \(\text{D. }y=13x-8\)
        \end{tabular*}\par}

        \item 已知抛物线 $C_1:y^2=2p_1x\,(p_1>0)$ 和 $C_2:x^2=2p_2y\,(p_2>0)$ 均经过点 $(4,8)$，则 $C_1$ 的焦点与 $C_2$ 的焦点之间的距离为

        {\centering\begin{tabular*}{\linewidth}{@{}@{\extracolsep{\fill}}llll@{}}
          \(\text{A. }12\)  & \(\text{B. }4\sqrt{5}\)  & \(\text{C. }6\)  & \(\text{D. }\dfrac{\sqrt{65}}{2}\)
        \end{tabular*}\par}

        \item 已知函数 $f(x)=\dfrac{x+2}{\e^x+a}$ 的最大值为 $1$，则 $a=$

        {\centering\begin{tabular*}{\linewidth}{@{}@{\extracolsep{\fill}}llll@{}}
          \(\text{A. }\dfrac{1}{2}\)  & \(\text{B. }1\)  & \(\text{C. }\dfrac{3}{2}\)  & \(\text{D. }2\)
        \end{tabular*}\par}

        \item 一百零八塔位于宁夏回族自治区青铜峡市，以其独特的建筑格局和深远的历史文化闻名遐迩. 该塔群共有 $108$ 座塔，依山势自上而下排成 $12$ 行，将第 $i$ 行中塔的座数记为 $a_i\,(i=1,2,\dots,12)$，已知 $a_1=1$，$a_2=a_3=3$，$a_4=a_5=5$，且 $a_6,a_7,\dots,a_{12}$ 是一个首项为 $7$，公差为 $2$ 的等差数列. 将 $a_1,a_2,\dots,a_{12}$ 分为 $6$ 组，每组 $2$ 个数，使得每组的 $2$ 个数之和可构成一个项数为 $6$ 且公差为 $d\,(d>0)$ 的等差数列，则 $d=$

        {\centering\begin{tabular*}{\linewidth}{@{}@{\extracolsep{\fill}}llll@{}}
          \(\text{A. }2\)  & \(\text{B. }4\)  & \(\text{C. }6\)  & \(\text{D. }8\)
        \end{tabular*}\par}

        \item 设 $U=\{(x_1,x_2,x_3)\mid x_i\in\{-2,-1,1,2\},i=1,2,3\}$ 为空间中 $64$ 个点构成的集合，点 $P(1,1,1)$，记样本空间 $\Omega=U\cup\{P\}$. 从 $\Omega$ 中随机取一个点，定义随机变量 $X$ 如下：对 $\Omega$ 中的每个点 $A(x_1,x_2,x_3)$，令 $X(A)=x_1+x_2+x_3$，则 $X$ 的数学期望为

        {\centering\begin{tabular*}{\linewidth}{@{}@{\extracolsep{\fill}}llll@{}}
          \(\text{A. }-\dfrac{1}{21}\)  & \(\text{B. }-\dfrac{1}{63}\)  & \(\text{C. }0\)  & \(\text{D. }\dfrac{1}{7}\)
        \end{tabular*}\par}
    \end{enumerate}
\end{description}

\begin{description}
    \item[二、选择题] 本题共3小题，每小题6分，共18分. 在每小题给出的选项中，有多项符合题目要求. 全部选对的得6分，部分选对的得部分分，有选错的得0分. 
    \begin{enumerate}[leftmargin=5pt]
      \addtocounter{enumi}{+8}   
        \item 设 $z=3+2\mathrm{i}$，则

        {\centering\begin{tabular*}{\linewidth}{@{}@{\extracolsep{\fill}}llll@{}}
          \(\text{A. }\bar{z}=3-2\mathrm{i}\)  & \(\text{B. }|z|=5\) &\(\text{C. }z^2=5+12\mathrm{i}\)  & \(\text{D. }\dfrac{z+3}{z-1}\in\mathbb{R}\)
        \end{tabular*}\par}

        \item 在空间中，$A,B$ 为两个定点，动点 $C$ 到直线 $AB$ 的距离为 $2$，动点 $D$ 到直线 $AB$ 的距离为 $1$，若二面角 $C-AB-D$ 为 $60^\circ$，则
        \begin{enumerate}[label=\Alph*.]
            \item $\angle CAD\geqslant 60^\circ$
            \item $CD\geqslant\sqrt{5}$
            \item 当 $AB\perp CD$ 时，$CD\perp$ 平面 $ABD$
            \item 当 $AB\perp$ 平面 $ACD$ 时，$AC\perp AD$
    \end{enumerate}

        \item 已知圆 $C_1:(x+1)^2+y^2=1$，圆 $C_2:(x-1)^2+y^2=1$，和圆 $C_3:x^2+(y-3)^2=1$，直线 $l:y=kx+b$ 与 $C_1,C_2,C_3$ 均有两个交点，记 $l$ 被 $C_1,C_2,C_3$ 截得的弦长分别为 $s_1,s_2,s_3$，则
        \begin{enumerate}[label=\Alph*.]
            \item $k$ 可以取任意实数
            \item 满足 $s_1=s_2=s_3$ 的直线 $l$ 共有 $3$ 条
            \item 满足 $s_1+s_2+s_3=3$ 的直线 $l$ 多于 $3$ 条
            \item 当 $b=0$ 时，$s_1+s_2+s_3$ 的最大值为 $\dfrac{2\sqrt{21}}{3}$
    \end{enumerate}
    \end{enumerate}
\end{description}

\begin{description}
    \item[三、填空题] 本题共3小题，每小题5分，共15分. 
    \begin{enumerate}[leftmargin=5pt]
              \addtocounter{enumi}{+11}   
        \item 双曲线 $3x^2-6y^2=1$ 的离心率等于 \rule[-0.3ex]{3em}{0.4pt}. 
        \item 已知 $f(x)=2\sin(ax+\theta)\,(a\in\mathbb{Z},0\leqslant\theta<2\pi)$ 是偶函数，$f(x)$ 在区间 $(0,\frac{\pi}{2})$ 单调递增，则 $a=$ \rule[-0.3ex]{3em}{0.4pt}，$f\bigl(\frac{2\pi}{3}\bigr)=$ \rule[-0.3ex]{3em}{0.4pt}. 
        \item 设实数 $q$ 满足：存在数列 $\{a_n\}$，使得对于任意 $n\in\mathbb{N}^*$，均有 $a_1+a_2+\cdots+a_{3n}=n^2+n$，且 $\{a_n\}$ 中有连续 $9$ 项 $a_k,a_{k+1},\dots,a_{k+8}$ 是公比为 $q$ 的等比数列. 则 $q$ 的最大值为 \rule[-0.3ex]{3em}{0.4pt}. 
    \end{enumerate}
\end{description}

\begin{description}
    \item[四、解答题] 本题共5小题，共77分. 解答应写出文字说明、证明过程或演算步骤. 
    \begin{enumerate}[leftmargin=5pt]
              \addtocounter{enumi}{+14}   
        \item （13分）

        如图，在直三棱柱 $ABC-A_1B_1C_1$ 中，$\angle ACB=90^\circ$，$AC=BC$，$D,E$ 分别为 $AB,AC_1$ 的中点. 

        \begin{enumerate}[label=(\arabic*)]
            \item 证明：$DE\parallel$ 平面 $BCC_1B_1$；
            \item 设 $CC_1=2$，直线 $DE$ 与平面 $ACC_1A_1$ 所成的角为 $45^\circ$，求直线 $DE$ 到平面 $BCC_1B_1$ 的距离. 
    \end{enumerate}

        \item （15分）

        已知在 $\triangle ABC$ 中，$AB=3$，$BC=2\sqrt{3}$，$\cos B=\dfrac{\sqrt{3}}{3}$. 

        \begin{enumerate}[label=(\arabic*)]
            \item 求 $\cos A$；
            \item 设 $D,E$ 两点满足：$D$ 在 $BA$ 的延长线上，$DE\parallel BC$，$AE\perp AC$. 若 $DE=6$，求 $CE$. 
    \end{enumerate}

        \item （15分）

        设整数 $N\geqslant 2$. 某同学用一个球进行投篮练习，至多投篮 $N$ 次. 当且仅当投中 $1$ 次时或 $N$ 次均未投中时，停止练习. 设该同学每次投中的概率为 $p\,(0<p<1)$，各次投中与否相互独立. 记 $X$ 为停止练习时该同学的投篮次数. 

\begin{enumerate}[label=(\arabic*)]
            \item 当 $N=4$，$p=\dfrac{1}{3}$ 时，求 $X$ 的分布列；
            \item 设 $k,m$ 均为自然数. 
            \begin{enumerate}[label=(\roman*)]
                \item 当 $k\leqslant N-1$ 时，求 $P(X>k)$；
                \item 当 $k+m\leqslant N-1$ 时，证明：$P(X>k+m\mid X>k)=P(X>m)$. 
\end{enumerate}
        \end{enumerate}

        \item （17分）

        已知椭圆 $C:\dfrac{x^2}{a^2}+\dfrac{y^2}{b^2}=1\,(a>b>0)$ 的左焦点为 $F(-1,0)$，离心率为 $\dfrac{1}{2}$. 

\begin{enumerate}[label=(\arabic*)]
            \item 求 $C$ 的方程；
            \item 设 $O$ 为坐标原点，过 $F$ 且斜率大于 $0$ 的动直线 $l$ 与 $C$ 交于 $P,Q$ 两点，其中 $P$ 在第三象限，直线 $PO$ 与 $C$ 的另一个交点为 $R$. 
            \begin{enumerate}[label=(\roman*)]
                \item 若 $\triangle PQR$ 的面积是 $\triangle PFO$ 的面积的 $3$ 倍，求 $l$ 的方程；
                \item 求 $\tan\angle PQR$ 的最小值. 
\end{enumerate}
        \end{enumerate}

        \item （17分）

        已知函数 $f(x)$ 的定义域为 $\mathbb{R}$，且当 $x<0$ 时，$f(x)=2x$，对任意 $x_0\in\mathbb{R}$，定义集合
        \[D(x_0)=\{d\in\mathbb{R}\mid f(x_0+d)>f(x_0)\}.\]

\begin{enumerate}[label=(\arabic*)]
            \item 若当 $x\geqslant 0$ 时，$f(x)=1-x$，求 $D(-1)$. 
            \item 若 $f(x)$ 是奇函数，$f(x_1)\leqslant f(x_2)$，且 $x_1x_2\neq 0$，证明：$D(x_2)\subseteq D(x_1)$；
            \item 若 $f(x)$ 满足：\ding{172} 若 $f(x_1)\leqslant f(x_2)$，则 $D(x_2)\subseteq D(x_1)$；\ding{173} 当 $0<x<1$ 时，$f(x)<f(0)$. 
            \begin{enumerate}[label=(\roman*)]
                \item 证明：$f(0)\geqslant 1$；
                \item 证明：$f(x)$ 在区间 $(0,+\infty)$ 单调递增. 
            \end{enumerate}
\end{enumerate}
    \end{enumerate}
\end{description}
\clearpage

\subsection{2026年全国II卷}\label{gk:2026年全国II卷}

\begin{description}
    \item[一、选择题] 本题共8小题，每小题5分，共40分．在每小题给出的四个选项中，只有一项是符合题目要求的．
    \begin{enumerate}[leftmargin=5pt]
        \item $(1-3\mathrm{i})^2=$
        
        {\centering\begin{tabular*}{\linewidth}{@{}@{\extracolsep{\fill}}llll@{}}
          \(\text{A. }-8+6\mathrm{i}\)  & \(\text{B. }-8-6\mathrm{i}\)  & \(\text{C. }8+6\mathrm{i}\)  & \(\text{D. }8-6\mathrm{i}\)
        \end{tabular*}\par}

        \item 已知向量 $\bm a,\bm b$ 满足 $|\bm a+\bm b|=1$，$|\bm a-\bm b|=3$，则 $\bm a\cdot\bm b=$
       
        {\centering\begin{tabular*}{\linewidth}{@{}@{\extracolsep{\fill}}llll@{}}
          \(\text{A. }\dfrac{1}{2}\)  & \(\text{B. }\dfrac{1}{4}\)  & \(\text{C. }-\dfrac{1}{2}\)  & \(\text{D. }-\dfrac{1}{4}\)
        \end{tabular*}\par}

        \item 已知集合 $A=\{0,1,3,6,9\}$，$B=\{x\mid x=\sqrt{x}\}$，则 $A\cap B=$
       
        {\centering\begin{tabular*}{\linewidth}{@{}@{\extracolsep{\fill}}llll@{}}
          \(\text{A. }\{0,1\}\)  & \(\text{B. }\{3,6\}\)  & \(\text{C. }\{0,1,9\}\)  & \(\text{D. }\{0,3,9\}\)
        \end{tabular*}\par}

        \item 双曲线 $C:\dfrac{x^2}{a^2}-\dfrac{y^2}{b^2}=1\,(a>0,b>0)$ 过点 $(1,0)$ 和 $(\dfrac{5}{2},3)$，则其渐近线方程为
        
        {\centering\begin{tabular*}{\linewidth}{@{}@{\extracolsep{\fill}}llll@{}}
          \(\text{A. }y=\pm3\sqrt{2}x\)  & \(\text{B. }y=\pm4\sqrt{3}x\)  & \(\text{C. }y=\pm\dfrac{3}{2}x\)  & \(\text{D. }y=\pm\dfrac{2}{6}x\)
        \end{tabular*}\par}

        \item 棱台上下底面均为有一个内角是 $60^\circ$ 的菱形，且上下底面边长分别为 $2$ 和 $3$，该棱台的高为 $3$，则该棱台体积为
      
        {\centering\begin{tabular*}{\linewidth}{@{}@{\extracolsep{\fill}}llll@{}}
          \(\text{A. }\dfrac{19}{12}\)  & \(\text{B. }\dfrac{19}{6}\)  & \(\text{C. }\dfrac{19}{4}\)  & \(\text{D. }\dfrac{19}{2}\)
        \end{tabular*}\par}

        \item 甲、乙、丙、丁等 $8$ 人分成 $A,B$ 两技术小组，要求每组 $4$ 人，且甲乙必须在同一组，丙丁不能在同一组，共有多少分配方案
       
        {\centering\begin{tabular*}{\linewidth}{@{}@{\extracolsep{\fill}}llll@{}}
          \(\text{A. }10\)  & \(\text{B. }12\)  & \(\text{C. }16\)  & \(\text{D. }24\)
        \end{tabular*}\par}

        \item 已知 $\alpha$ 为第二象限角，且 $3\sin2\alpha\cos\alpha=8\sin\alpha\cos2\alpha$，则 $\dfrac{1+\sin\alpha}{2-\cos\alpha}=$
       
        {\centering\begin{tabular*}{\linewidth}{@{}@{\extracolsep{\fill}}llll@{}}
          \(\text{A. }\dfrac{3}{4}\)  & \(\text{B. }\dfrac{3}{2}\)  & \(\text{C. }\dfrac{1}{2}\)  & \(\text{D. }\dfrac{15}{8}\)
        \end{tabular*}\par}

        \item 已知 $f(x)$ 为定义在 $\mathbb{R}$ 上的偶函数，且 $f(x)+f(x-2)=0$，当 $x\in\bigl[\dfrac{3}{2},3\bigr]$ 时，$f(x)=x^2+ax+b$，则
        
        {\centering\begin{tabular*}{\linewidth}{@{}@{\extracolsep{\fill}}llll@{}}
          \(\text{A. }a=-2,b=-3\)  & \(\text{B. }a=-2,b=3\)  & \(\text{C. }a=-4,b=-3\)  & \(\text{D. }a=-4,b=3\)
        \end{tabular*}\par}
    \end{enumerate}
\end{description}

\begin{description}
    \item[二、选择题] 本题共3小题，每小题6分，共18分．在每小题给出的选项中，有多项符合题目要求．全部选对的得6分，部分选对的得部分分，有选错的得0分．
    \begin{enumerate}[leftmargin=5pt]
      \addtocounter{enumi}{+8}
        \item 已知 $\odot O:x^2+y^2=1$，$\odot A:x^2+y^2-6x-8y+k=0$，则
        \begin{enumerate}[label=\Alph*.]
            \item 点 $A$ 的坐标为 $(-3,-4)$
            \item $k=9$ 时，$\odot A$ 与 $x$ 轴相切
            \item 当 $k=-11$ 时，$\odot A$ 与 $\odot O$ 相切
            \item 当 $\odot O$ 与 $\odot A$ 相交时，两交点所在直线的方程是 $6x+8y-k-2=0$
        \end{enumerate}

        \item 等比数列 $\{a_n\}$ 的公比 $q\neq1$，$a_1>0$，$2a_3=a_2+a_1$，记前 $n$ 项和为 $S_n$，则
        \begin{enumerate}[label=\Alph*.]
            \item $q=-\dfrac{1}{2}$
            \item $S_n>\dfrac{2a_1}{3}$
            \item $2S_{n+2}=S_{n+1}+S_n$
            \item $\displaystyle\sum_{k=1}^{n}S_k>\dfrac{2na_1}{3}$
        \end{enumerate}

        \item 已知抛物线 $E:y^2=8x$，斜率 $k\,(k>0)$ 的直线 $l$ 过点 $(1,0)$，$\triangle ABC$ 为等边三角形，$A$ 在 $y$ 轴上，$B,C$ 在 $l$ 上，则
        \begin{enumerate}[label=\Alph*.]
            \item 抛物线准线方程为 $x=-2$
            \item $l$ 与 $y$ 轴交点为 $(0,-k)$
            \item 若 $l$ 与 $E$ 相交于唯一点 $B$，则抛物线焦点在直线 $AB$ 上
            \item $k=2$ 时，$\triangle ABC$ 面积最小值为 $\dfrac{3}{2}$
        \end{enumerate}
    \end{enumerate}
\end{description}

\begin{description}
    \item[三、填空题] 本题共3小题，每小题5分，共15分．
    \begin{enumerate}[leftmargin=5pt]
      \addtocounter{enumi}{+11}
        \item $S_n$ 为等差数列 $\{a_n\}$ 前 $n$ 项和．若 $a_1=-1$，$a_4=5$，则 $S_6=$ \rule[-0.3ex]{3em}{0.4pt}．
        \item 若函数 $f(x)=2^x+2^{-x}-m$ 有两个零点，则 $m$ 的取值范围是 \rule[-0.3ex]{3em}{0.4pt}．
        \item 球 $O$ 的体积为 $4\sqrt{3}\pi$，$A,B,C,D$ 四点均在球 $O$ 的球面上，$\triangle ABC$ 为等边三角形，$DA=DB=DC=2$，则 $\triangle ABC$ 的面积为 \rule[-0.3ex]{3em}{0.4pt}．
    \end{enumerate}
\end{description}

\begin{description}
    \item[四、解答题] 本题共5小题，共77分．解答应写出文字说明、证明过程或演算步骤．
    \begin{enumerate}[leftmargin=5pt]
      \addtocounter{enumi}{+14}
        \item （13分）

        某工厂抽取一批电子元件检测，记录第一次出现故障的时间（天），绘制成如下的频率分布直方图：

        \begin{enumerate}[label=(\arabic*)]
            \item 求第一四分位数和中位数；
            \item $\hat{p}$ 为首次故障时间小于 $365$ 天的概率估计值．
            \begin{enumerate}[label=(\roman*)]
                \item 求 $\hat{p}$；
                \item 工厂向某用户销售 $100$ 件电子元件，$X$ 为这 $100$ 件产品首次出现故障小于 $365$ 天的件数，则 $X\sim \mathrm{B}(100,\hat{p})$，求 $\mathrm{E}(X)$，$\mathrm{D}(X)$．
            \end{enumerate}
        \end{enumerate}

        \item （15分）

        三棱锥 $A-BCD$ 中，$E$ 在 $BD$ 上，$AE\perp CE$，$AE\perp DE$，$CD\perp AD$．

        \begin{enumerate}[label=(\arabic*)]
            \item 证明：$CD\perp AB$；
            \item 若 $DE=2$，$BE=1$，$AE=2$，$CD=2\sqrt{3}$，求 $AD$ 与平面 $ABC$ 所成角的正弦值．
        \end{enumerate}

        \item （15分）

        在 $\triangle ABC$ 中，已知 $\cos B=\dfrac{3}{4}$，$\cos2(A+C)+\sin A\sin C=1$．

        \begin{enumerate}[label=(\arabic*)]
            \item 证明：$\triangle ABC$ 为钝角三角形；
            \item 若 $\triangle ABC$ 面积为 $\dfrac{7}{4}$，求 $\triangle ABC$ 周长．
        \end{enumerate}

        \item （17分）

        椭圆 $E:\dfrac{x^2}{a^2}+y^2=1\,(a>1)$，过右焦点垂直于 $x$ 轴的直线被 $E$ 所截线段长为 $2$．

        \begin{enumerate}[label=(\arabic*)]
            \item 求 $E$ 的离心率；
            \item $O$ 为坐标原点，给定点 $G(t_0,0)\,(t_0\neq0)$；$A(x_0,y_0)\,(y_0\neq0)$ 在 $E$ 上，过点 $A$ 作 $y$ 轴的垂线，交于点 $B$，$AO$ 与 $GB$ 交于点 $P$．当 $A$ 在 $E$ 上运动时，$P$ 的轨迹为 $M$．
            \begin{enumerate}[label=(\roman*)]
                \item 求 $M$ 的方程；
                \item $M$ 是否有中心点？当 $t_0$ 为何值时，$M$ 有中心点？当 $M$ 有中心点时，平移 $M$ 到 $M'$，使 $O$ 为 $M'$ 的中心点，说明 $M'$ 为何形状？
            \end{enumerate}
        \end{enumerate}

        \item （17分）

        已知函数 $f(x)=x\e^x+ax+b$，曲线 $y=f(x)$ 在点 $(0,f(0))$ 处的切线方程为 $y=-2x+1$．

        \begin{enumerate}[label=(\arabic*)]
            \item 求 $a,b$；
            \item 当 $x>0$ 时，$f(x+m)-f(x)>m$，求 $m$ 的取值范围；
            \item 当 $x>0$ 时，$f(x+k)+f(k-x)>2f(k)$，求 $k$ 的最小值．
        \end{enumerate}
    \end{enumerate}
\end{description}
\clearpage

\subsection{2026年北京卷}\label{gk:2026年北京卷}

\begin{description}
    \item[一、选择题] 本题共10小题，每小题4分，共40分．在每小题给出的四个选项中，选出符合题目要求的一项．
    \begin{enumerate}[leftmargin=5pt]
        \item 已知集合 $M=\{x\mid -1<x<3\}$，$N=\{x\mid x\geqslant 2\}$，则 $M\cup N=$
       
        {\centering\begin{tabular*}{\linewidth}{@{}@{\extracolsep{\fill}}llll@{}}
          \(\text{A. }(-1,+\infty)\)  & \(\text{B. }[2,3)\)  & \(\text{C. }(-1,3]\)  & \(\text{D. }\mathbb{R}\)
        \end{tabular*}\par}

        \item 已知复数 $z_1,z_2$，则 $|z_1+z_2|=$
       
        {\centering\begin{tabular*}{\linewidth}{@{}@{\extracolsep{\fill}}llll@{}}
          \(\text{A. }1\)  & \(\text{B. }\sqrt{2}\)  & \(\text{C. }\sqrt{3}\)  & \(\text{D. }2\)
        \end{tabular*}\par}

        \item 已知双曲线 $\dfrac{x^2}{a^2}-\dfrac{y^2}{4}=1\,(a>0)$ 的一条渐近线为 $y=\dfrac{2}{3}x$，则 $a=$
       
        {\centering\begin{tabular*}{\linewidth}{@{}@{\extracolsep{\fill}}llll@{}}
          \(\text{A. }2\)  & \(\text{B. }3\)  & \(\text{C. }4\)  & \(\text{D. }9\)
        \end{tabular*}\par}

        \item 在 $(a-\sqrt{x})^7$ 的展开式中，$x^2$ 的系数为 $280$，则 $a=$
       
        {\centering\begin{tabular*}{\linewidth}{@{}@{\extracolsep{\fill}}llll@{}}
          \(\text{A. }2\)  & \(\text{B. }-2\)  & \(\text{C. }\sqrt{2}\)  & \(\text{D. }\pm2\)
        \end{tabular*}\par}

        \item 下列函数中，既是奇函数，又在其定义域内单调递增的是
       
        {\centering\begin{tabular*}{\linewidth}{@{}@{\extracolsep{\fill}}llll@{}}
          \(\text{A. }f(x)=x^2\)  & \(\text{B. }f(x)=\sin x\)  & \(\text{C. }f(x)=-x^3\)  & \(\text{D. }f(x)=\ln\dfrac{5+x}{5-x}\)
        \end{tabular*}\par}

        \item 已知向量 $\bm a,\bm b$ 满足 $\bm a=(2,0)$，$|\bm a-\bm b|=2$，则 $|\bm b|$ 的最大值为
       
        {\centering\begin{tabular*}{\linewidth}{@{}@{\extracolsep{\fill}}llll@{}}
          \(\text{A. }1\)  & \(\text{B. }2\)  & \(\text{C. }3\)  & \(\text{D. }4\)
        \end{tabular*}\par}

        \item 已知数列 $\{a_n\},\{b_n\}$．
        
        命题 $P$：存在常数 $M$，使 $a_n\leqslant M\leqslant b_n$ 对一切 $n$ 成立；
        
        命题 $Q$：$a_n\leqslant b_n$ 对一切 $n$ 成立．
        
        则 $P$ 是 $Q$ 的
       
        {\centering\begin{tabular*}{\linewidth}{@{}@{\extracolsep{\fill}}llll@{}}
          \(\text{A. }\)充分不必要条件  & \(\text{B. }\)必要不充分条件  & \(\text{C. }\)充要条件  & \(\text{D. }\)既不充分也不必要条件
        \end{tabular*}\par}

        \item 已知 $f(x)=\sin(x+\varphi)\,(0\leqslant\varphi<2\pi)$．将 $f(x)$ 的图象向右平移 $3\varphi$ 个单位得 $g(x)$，若 $g(x)$ 的图象与 $f(x)$ 的图象关于 $x$ 轴对称，则满足条件的 $\varphi$ 的个数为
       
        {\centering\begin{tabular*}{\linewidth}{@{}@{\extracolsep{\fill}}llll@{}}
          \(\text{A. }1\)  & \(\text{B. }2\)  & \(\text{C. }3\)  & \(\text{D. }4\)
        \end{tabular*}\par}

        \item 某校组织高一、高二年级学生分别前往甲、乙两地研学．记高一去甲、乙地的人数为 $p,q$，高二去甲、乙地的人数为 $r,s$．已知 $p+q>r+s$ 且 $p+r>q+s$．下列关于人数的判断中一定正确的是
       
        {\centering\begin{tabular*}{\linewidth}{@{}@{\extracolsep{\fill}}llll@{}}
          \(\text{A. }\)去甲的高一学生多于去乙的高二学生 \\
          \(\text{B. }\)去甲的高一学生不多于去乙的高二学生 \\
          \(\text{C. }\)去乙的高一学生多于去甲的高二学生 \\
          \(\text{D. }\)以上都不能确定
        \end{tabular*}\par}

        \item 已知 $A,B$ 是平面上两个定点且 $AB=4$，点 $D$ 在以 $A$ 为圆心、半径 $1$ 的圆上，点 $C$ 满足 $CD=3$，$BC=\dfrac{5}{2}$．则 $\cos\angle ABC$ 的取值范围为
       
        {\centering\begin{tabular*}{\linewidth}{@{}@{\extracolsep{\fill}}llll@{}}
          \(\text{A. }\left[\dfrac{1}{4},\dfrac{3}{4}\right]\)  & \(\text{B. }\left[\dfrac{1}{2},\dfrac{3}{4}\right]\)  & \(\text{C. }\left[\dfrac{1}{4},\dfrac{1}{2}\right]\)  & \(\text{D. }\left[\dfrac{1}{2},1\right]\)
        \end{tabular*}\par}
    \end{enumerate}
\end{description}

\begin{description}
    \item[二、填空题] 本题共5小题，每小题5分，共25分．
    \begin{enumerate}[leftmargin=5pt]
      \addtocounter{enumi}{+10}
        \item 若直线 $ax+y=0$ 与圆 $(x-2)^2+(y-2)^2=4$ 相切，则 $a=$ \rule[-0.3ex]{3em}{0.4pt}．
        \item 等差数列 $\{a_n\}$ 的前 $n$ 项和为 $S_n$，满足 $S_6=6a_4+30$，且对任意 $n$ 有 $S_n\leqslant S_5$，则 $a_5$ 的一个可能取值为 \rule[-0.3ex]{3em}{0.4pt}．
        \item 声压级 $y$（单位：$\text{dB}$）与频率 $f$（单位：$\text{Hz}$）满足 $y=\lg\bigl(1+\frac{f}{700}\bigr)$．若 $\lg2\leqslant y<3\lg2$，则 $f$ 的取值范围为 \rule[-0.3ex]{3em}{0.4pt}．
        \item 三棱锥 $A$-$BCD$ 中，$AD=AB=AC=2\sqrt{2}$，$BD=BC=2$，$DC=2\sqrt{3}$，则其底面 $BCD$ 的面积为 \rule[-0.3ex]{3em}{0.4pt}，体积为 \rule[-0.3ex]{3em}{0.4pt}．
        \item 设 $c\in\mathbb{R}$，函数 $f(x)=|x^2-\cos x-c^2|$．给出下列四个命题：
        \begin{enumerate}[label=\textcircled{\arabic*}]
            \item $f(x)$ 在 $(-1,1]$ 上既有最小值又有最大值；
            \item 当 $c=0$ 时，$f(x)=1$ 有 $3$ 个解；
            \item 当 $c=1$，$x\in(1,2]$ 时，$f(x)$ 有最大值；
            \item 当 $c>0$ 时，$f(x)$ 与 $y=c$ 有 $4$ 个交点．
        \end{enumerate}
        其中正确命题的序号是 \rule[-0.3ex]{3em}{0.4pt}．
    \end{enumerate}
\end{description}

\begin{description}
    \item[三、解答题] 本题共6小题，共85分．解答应写出文字说明、证明过程或演算步骤．
    \begin{enumerate}[leftmargin=5pt]
      \addtocounter{enumi}{+15}
        \item （13分）

        已知函数 $f(x)=2\sin\omega x\cos\varphi+2\cos\omega x\sin\varphi\,(\omega>0,\,0<\varphi<\frac{\pi}{2})$ 的最小正周期为 $\pi$，且 $f\bigl(\frac{\pi}{4}\bigr)=1$．

        \begin{enumerate}[label=(\arabic*)]
            \item 求 $\omega,\varphi$ 的值；
            \item 求 $f(x)$ 的单调递减区间．
        \end{enumerate}

        \item （14分）

        现从全校学生中随机抽取 $200$ 人统计数学成绩，成绩分组及对应人数如下：

        {\centering\begin{tabular}{|c|c|c|c|c|c|}
        \hline
        成绩分组 & $[81,94]$ & $(94,107]$ & $(107,120]$ & $(120,135]$ & $(135,150]$ \\
        \hline
        人数 & $40$ & $60$ & $60$ & $28$ & $12$ \\
        \hline
        \end{tabular}\par}

        以频率估计概率，完成下列问题．

        \begin{enumerate}[label=(\arabic*)]
            \item 求数学成绩不高于 $120$ 分的概率；
            \item 从全校学生中随机抽取 $4$ 人，求其中恰好 $2$ 人数学成绩不低于 $120$ 且 $2$ 人数学成绩不高于 $94$ 的概率；
            \item 若把每组数据分别用其区间的左端点、中点、右端点代表，所得三组数据的方差分别记为 $s_1^2,s_2^2,s_3^2$，试比较其大小并说明理由．
        \end{enumerate}

        \item （14分）

        如图，直三棱柱 $ABC$-$A_1B_1C_1$ 中，$\angle BAC=90^\circ$，$AB=AC=2$，$BB_1=\sqrt{2}$，$E,D$ 分别为 $A_1B_1,AC$ 的中点．

        \textsf{[插入图片：直三棱柱示意图]}

        \begin{enumerate}[label=(\arabic*)]
            \item 求证：$DE\parallel$ 平面 $BB_1C_1C$；
            \item 点 $P$ 在平面 $A_1B_1C_1$ 内，且 $EP\parallel B_1C_1$，再从条件一、条件二、条件三中选择一个作为已知，使得 $P$ 唯一确定，求平面 $PAD$ 与平面 $PDE$ 的夹角的余弦值．
            \begin{quote}
            条件一：$PA=PD$；条件二：$PA\perp BC$；条件三：$BB_1\parallel$ 平面 $PDE$．
            \end{quote}
        \end{enumerate}

        \item （15分）

        已知椭圆 $E:\dfrac{x^2}{a^2}+\dfrac{y^2}{b^2}=1\,(a>b>0)$ 的一个顶点是 $(2,0)$，离心率为 $\dfrac{1}{2}$．

        \begin{enumerate}[label=(\arabic*)]
            \item 求 $E$ 的方程；
            \item 过点 $A(1,1)$ 作斜率为 $k\,(k\neq\pm1)$ 的直线交 $E$ 于 $B,C$ 两点．设 $D$ 为 $B$ 关于直线 $y=x$ 的对称点，直线 $DC$ 交 $y=x$ 于点 $Q$．若 $\triangle ABD$ 与 $\triangle ACQ$ 的面积相差 $\dfrac{5}{8}$，求 $k$ 的值．
        \end{enumerate}

        \item （15分）

        已知函数 $f(x)=x^2+mx-4\e^{nx-1}$，其中 $n\in\mathbb{Z}$．曲线 $y=f(x)$ 在点 $(-1,f(-1))$ 处的切线方程为 $4x+\e^2y+\e^2+8=0$．

        \begin{enumerate}[label=(\arabic*)]
            \item 求 $m,n$ 的值；
            \item 求证：$f(x)$ 有两个极值点；
            \item 当 $k>0$ 时，讨论直线 $y=kx-1$ 与曲线 $y=f(x)$ 的公共点个数．
        \end{enumerate}

        \item （15分）

        已知 $A$ 是 $m\times n$ 的数表
        \[
        \begin{bmatrix}
        a_{11}&a_{12}&\cdots&a_{1n}\\
        a_{21}&a_{22}&\cdots&a_{2n}\\
        \vdots&\vdots&\ddots&\vdots\\
        a_{m1}&a_{m2}&\cdots&a_{mn}
        \end{bmatrix}
        \]
        其中 $a_{ij}\in\{-1,1\}$（$1\leqslant i\leqslant m$，$1\leqslant j\leqslant n$，且 $m,n\in\mathbb{N}^*$）．若对任意行标 $l,k$（$l>k$）以及列标 $p,q$（$p>q$），满足 $(l-k)-(p-q)\in\{-2,2\}$，则有
        \[
        a_{lp}+a_{lq}+a_{kp}+a_{kq}=0,
        \]
        称 $A$ 具有性质 $P$．

        \begin{enumerate}[label=(\arabic*)]
            \item 已知数表
            \[
            A_1:\begin{bmatrix}1&1&-1&1\\ -1&1&-1&1\end{bmatrix}\quad
            A_2:\begin{bmatrix}1&-1&1&-1\\ 1&1&-1&-1\\ -1&1&1&1\end{bmatrix}
            \]
            判断 $A_1,A_2$ 是否具有性质 $P$；
            \item 已知 $4\times5$ 的数表 $A$ 具有性质 $P$，求其中取值为 $1$ 的元素个数的最大值；
            \item 已知 $6\times6$ 的数表 $A$ 具有性质 $P$，且 $a_{11}=1$，求证：对任意 $i,j\in\mathbb{N}^*$ 且 $1\leqslant i,j\leqslant6$，均有 $a_{ij}=a_{i1}\cdot a_{1j}$．
        \end{enumerate}
    \end{enumerate}
\end{description}
\clearpage


\subsection{2026年上海卷}\label{gk:2026年上海卷}

\begin{description}
    \item[一、填空题] 本大题共有12题，满分54分，第1\textasciitilde6题每题4分，第7\textasciitilde12题每题5分．
    \begin{enumerate}[leftmargin=5pt]
        \item 已知集合 $A=\{2,1+a\}$，$-1\in A$，则 $a=$ \rule[-0.3ex]{3em}{0.4pt}．
        \item 已知 $\{a_n\}$ 为等比数列，$a_1=2$，$a_2=6$，则 $a_4=$ \rule[-0.3ex]{3em}{0.4pt}．
        \item 已知 $\sin\alpha=\dfrac{1}{5}$，则 $\cos2\alpha=$ \rule[-0.3ex]{3em}{0.4pt}．
        \item 已知事件 $A,B$ 互斥，$P(A)=0.2$，$P(B)=0.5$，则 $P(A\cup B)=$ \rule[-0.3ex]{3em}{0.4pt}．
        \item 已知函数 $f(x)$ 是偶函数，当 $x\geqslant0$ 时，$f(x)=\sqrt{x}-a$，若 $f(-4)=3$，则 $a=$ \rule[-0.3ex]{3em}{0.4pt}．
        \item 已知 $(x^2+x)^5$，则展开式中 $x^7$ 的系数为 \rule[-0.3ex]{3em}{0.4pt}．
        \item 已知 $a^2+4b^2=1$，则 $ab$ 的最大值为 \rule[-0.3ex]{3em}{0.4pt}．
        \item 已知随机变量 $X$ 的分布为 $\begin{pmatrix}-1 & 0 & 1 \\ a & 0.3 & b\end{pmatrix}$ 且 $E(X)=0.5$，则 $b=$ \rule[-0.3ex]{3em}{0.4pt}．
        \item 已知等差数列 $\{a_n\}$ 中，$a_1=0$，公差为 $d$，其前 $n$ 项和 $S_n$ 在区间 $(0,1)$ 内至少有两项，则公差 $d$ 的取值范围是 \rule[-0.3ex]{3em}{0.4pt}．
        \item 已知向量 $\vec{a},\vec{b},\vec{c}$ 是两两不平行的向量，且 $\vec{a}+3\vec{b}\parallel\vec{b}+\vec{c}$，$2\vec{a}+k\vec{c}\parallel\vec{b}+\vec{c}$，则 $k$ 的值为 \rule[-0.3ex]{3em}{0.4pt}．
        \item 已知三角函数 $f(t)=A\sin(\omega t+\varphi)+B\,(A>0,B\in\mathbb{R},\omega>0,0\leqslant\varphi<2\pi)$，若 $v=f(t)$，当 $v=0$ 或 $v=4$ 时其导数为 $0$，初始速度为 $0$，且速度第一次达到 $4$ 时用时为 $0.1$ 秒，则 $f(t)=$ \rule[-0.3ex]{3em}{0.4pt}．
        \item 在 $\triangle ABC$ 中，$AB=3$，$AC=5$，$BC=\sqrt{14}$．已知点 $A,B,C$ 分别为椭圆的上、下、右顶点以及两个焦点中的三点，求椭圆的离心率． \rule[-0.3ex]{3em}{0.4pt}．
    \end{enumerate}
\end{description}

\begin{description}
    \item[二、选择题] 本题共有4题，满分18分，13、14每题4分，15、16题每题5分．每题有且只有一个正确选项．
    \begin{enumerate}[leftmargin=5pt]
      \addtocounter{enumi}{+12}
        \item $a$ 为不为 $1$ 的任意实数，则 $a\cdot\sqrt[3]{a}=$
       
        {\centering\begin{tabular*}{\linewidth}{@{}@{\extracolsep{\fill}}llll@{}}
          \(\text{A. }a^{\frac{3}{2}}\) & \(\text{B. }a^{\frac{4}{3}}\) & \(\text{C. }a^{\frac{5}{2}}\) & \(\text{D. }a^{\frac{5}{3}}\)
        \end{tabular*}\par}

        \item 事件 $A$ 和事件 $B$ 相互独立，``$A$ 和 $B$ 至少一个发生''的对立事件是
       
        {\centering\begin{tabular*}{\linewidth}{@{}@{\extracolsep{\fill}}llll@{}}
          \(\text{A. }A\cap B\) & \(\text{B. }A\cup B\) & \(\text{C. }\overline{A}\cap\overline{B}\) & \(\text{D. }\overline{A}\cup\overline{B}\)
        \end{tabular*}\par}

        \item 已知 $w,z$ 为复数，当 $w-z$ 为实数或 $w-z$ 的共轭复数为实数时，称 $w$ 和 $z$ 互相伴随．则当 $w$ 和 $z$ 互相伴随时，$w+\mathrm{i}$ 和 $z+\mathrm{i}$ 互相伴随的充要条件是
       
        {\centering\begin{tabular*}{\linewidth}{@{}@{\extracolsep{\fill}}llll@{}}
          \(\text{A. }\operatorname{Re}w+\operatorname{Re}z=0\) & \(\text{B. }\operatorname{Re}w-\operatorname{Re}z=0\) & \(\text{C. }\operatorname{Im}w+\operatorname{Im}z=0\) & \(\text{D. }\operatorname{Im}w-\operatorname{Im}z=0\)
        \end{tabular*}\par}

        \item 已知空间直角坐标系中有一正方体，其三组棱分别与 $x$ 轴、$y$ 轴、$z$ 轴重合，顶点 $A$ 与坐标原点重合，点 $C$ 是正方体底面中与 $A$ 相对的对角顶点，点 $C_1$ 在点 $C$ 的正上方．将正方体绕直线 $AC_1$ 旋转一周，试问点 $C$ 的运动轨迹会经过几个空间卦限
       
        {\centering\begin{tabular*}{\linewidth}{@{}@{\extracolsep{\fill}}llll@{}}
          \(\text{A. }1\) & \(\text{B. }3\) & \(\text{C. }4\) & \(\text{D. }7\)
        \end{tabular*}\par}
    \end{enumerate}
\end{description}

\begin{description}
    \item[三、解答题] 本大题共有5题，满分78分．解答下列各题必须写出必要的步骤．
    \begin{enumerate}[leftmargin=5pt]
      \addtocounter{enumi}{+16}
        \item （14分）

        某一年某一样物质，记录了颗粒物密度和二氧化硫密度，下表为2023年前九年数据：
        \begin{table}[H]
        \centering
        \begin{tabular}{|c|c|c|c|c|c|c|c|c|c|}
        \hline 颗粒物密度 & 101.02 & 87.02 & 57.46 & 21.85 & 11.76 & 8.86 & 5.03 & 4.63 & 3.86 \\
        \hline 二氧化硫密度 & 119.47 & 81.94 & 53.20 & 9.16 & 6.60 & 4.40 & 3.31 & 3.35 & 3.86 \\
        \hline
        \end{tabular}
        \end{table}

        \begin{enumerate}[label=(\arabic*)]
            \item 从九年中任意抽取一年，求颗粒物密度比二氧化硫密度高的概率；
            \item 想要研究颗粒物密度和二氧化硫密度的线性关系，应选择扇形图、茎叶图还是散点图？二者的相关系数落在区间 $(-1,0)$、$(0,+\infty)$ 还是 $(1,2)$？直接写出答案即可；
            \item 已知年份的平均数是 $2018$，现有两个拟合函数：$y_1=106.55e^{-0.461x}$，$y_2=a(x-2014)+83.57\,(a\in\mathbb{R})$，判断哪一个函数的预测值与实际值之差的绝对值更小（精确到 $0.01$），并说明理由．
        \end{enumerate}

        \item （14分）

        已知四棱锥 $P-ABCD$，底面 $ABCD$ 为矩形，$PH\perp$ 底面 $ABCD$，垂足 $H$ 在 $CD$ 边上，且 $CH=1$，$HD=4$，$BC=2$．

        \begin{enumerate}[label=(\arabic*)]
            \item 求证：$AH\perp PB$；
            \item 若四棱锥 $P-ABCD$ 的体积为 $\dfrac{10\sqrt{5}}{3}$，求二面角 $C-HB-P$ 的大小．
        \end{enumerate}

        \item （14分）

        已知函数 $f(x)=x^2+ax+3$，$g(x)=4x+\dfrac{1}{x^2}$．

        \begin{enumerate}[label=(\arabic*)]
            \item 解不等式：$f(x)+\dfrac{1}{x^2}>g(x)$；
            \item 设直线 $l_1$ 为 $f(x)$ 过点 $(0,3)$ 的切线，直线 $l_2\perp l_1$ 且也过点 $(0,3)$．若 $l_1,l_2$ 与 $g(x)$ 在第一象限内无交点，求实数 $a$ 的取值范围．
        \end{enumerate}

        \item （18分）

        已知双曲线 $\Gamma:x^2-y^2=1$，点 $P$ 在 $\Gamma$ 上，$F_1,F_2$ 分别为双曲线的左、右焦点．

        \begin{enumerate}[label=(\arabic*)]
            \item 求点 $(2,0)$ 到双曲线渐近线的距离；
            \item 若 $\overrightarrow{PF_1}\cdot\overrightarrow{PF_2}=1$，求 $S_{\triangle PF_1F_2}$；
            \item 记 $\Omega$ 为双曲线 $\Gamma$ 满足 $\begin{cases}x>0 \\ y\geqslant-1\end{cases}$ 与 $\begin{cases}x<0 \\ y\leqslant-1\end{cases}$ 的部分；直线 $l,m$ 均过右焦点 $F_2$，$l$ 与 $\Omega$ 交于 $P,Q$ 两点（分别在第一、第四象限），$m$ 与 $\Omega$ 交于 $M,N$ 两点（分别在第三、第四象限）．问：是否存在常数 $\lambda$，使得对任意直线 $l$，都存在唯一对应的直线 $m$ 满足 $\lambda|PQ|=|MN|$？若存在，求出 $\lambda$ 的值；若不存在，请说明理由．
        \end{enumerate}

        \item （18分）

        已知 $(i,j,k)$ 是 $1,2,3$ 的一个排列，若函数 $f_1(x),f_2(x),f_3(x)$ 满足：对任意 $x\in I$，都有 $f_i(x)\leqslant f_j(x)$ 且 $f_i(x)+f_j(x)\leqslant f_k(x)+f_j(x)$，则称 $(i,j,k)$ 是关于 $f_1(x),f_2(x),f_3(x)$ 的一个 $I$ 排列，关于 $f_1(x),f_2(x),f_3(x)$ 的 $I$ 排列总数记为 $n_I$．

        \begin{enumerate}[label=(\arabic*)]
            \item 已知 $I=[3,+\infty)$，$f_1(x)=x$，$f_2(x)=0$，$f_3(x)=x^2+1$，判断 $(3,1,2)$ 是否为 $I$ 排列；
            \item 设 $I=(0,+\infty)$，$f_1(x)=x-1$，$f_2(x)=x+m$，$f_3(x)=x^2$，若满足条件的 $n_I=6$，求 $m$ 的取值范围；
            \item 设 $x\in[0,+\infty)$ 时恒有 $0<F(x)<1$，令 $I=[a,+\infty)$，$f_1(x)=F(x)$，$f_2(x)=\dfrac{1}{2}\bigl(F(x+a)+F(x-a)\bigr)$，$f_3(x)=1-e^{-x}$．证明：若 $F(x)$ 严格递减，则存在 $a>0$，使 $n_I\geqslant4$；若 $F(x)$ 严格递增，则存在 $a\in(0,1)$，使得 $n_I\neq2$．
        \end{enumerate}
    \end{enumerate}
\end{description}
\clearpage

\subsection{2026年天津卷}\label{gk:2026年天津卷}

\begin{description}
    \item[一、选择题] 本题共9小题，每小题5分，共45分．在每小题给出的四个选项中，只有一项是符合题目要求的．
    \begin{enumerate}[leftmargin=5pt]
        \item 已知全集 $U=\{-2,-1,0,1,2,3\}$，集合 $A=\{-1,0,1,3\}$，集合 $B=\{-2,0,1\}$，则 $(\complement_U A)\cup B=$
        
        {\centering\begin{tabular*}{\linewidth}{@{}@{\extracolsep{\fill}}llll@{}}
          \(\text{A. }\{-2\}\) & \(\text{B. }\{-2,2\}\) & \(\text{C. }\{0,1,2\}\) & \(\text{D. }\{-2,0,1,2\}\)
        \end{tabular*}\par}

        \item 设 $x\in\mathbb{R}$，则``$x>0$''是``$x^2+3x>0$''的
        
        {\centering\begin{tabular*}{\linewidth}{@{}@{\extracolsep{\fill}}llll@{}}
          \(\text{A. }\)充分而不必要条件 & \(\text{B. }\)必要而不充分条件 & \(\text{C. }\)充分必要条件 & \(\text{D. }\)既不充分也不必要条件
        \end{tabular*}\par}

        \item 调查候鸟和温度的关系，在不同温度下统计候鸟的数量，所得数据如图所示，其中相关系数 $r=-0.91$，根据最小二乘法算得：$\hat{y}=-1.17x+1370.7$，下列说法正确的是
        
        {\centering\begin{tabular*}{\linewidth}{@{}@{\extracolsep{\fill}}llll@{}}
          \(\text{A. }y\text{与}x\text{负相关}\) & \(\text{B. 当}x=10\text{时, }y\text{一定为}1359\) & \(\text{C. 当}x=10\text{时, }y\text{一定小于}1359\) & \(\text{D. }\)两变量无线性关系
        \end{tabular*}\par}

        \item 函数 $f(x)$ 的部分图象如图所示，则 $f(x)$ 的解析式可能为
        
        {\centering\begin{tabular*}{\linewidth}{@{}@{\extracolsep{\fill}}llll@{}}
          \(\text{A. }x+\sin\pi x\) & \(\text{B. }x-\sin\pi x\) & \(\text{C. }-x+\sin\pi x\) & \(\text{D. }x\cdot\sin\pi x\)
        \end{tabular*}\par}

        \item 正方体 $ABCD-A_1B_1C_1D_1$ 中，错误的是
       
        {\centering\begin{tabular*}{\linewidth}{@{}@{\extracolsep{\fill}}llll@{}}
          \(\text{A. }AC\parallel A_1C_1\) & \(\text{B. }CC_1\perp\text{面}ABC\) & \(\text{C. }ACD_1\parallel\text{面}A_1C_1B\) & \(\text{D. }ADD_1\perp\text{面}ACD_1\)
        \end{tabular*}\par}

        \item 已知函数 $f(x)=|\ln x|$．若 $a=f(2^{0.3})$，$b=f(3^{0.3})$，$c=f(3^{-0.5})$，则 $a,b,c$ 的大小关系为
        
        {\centering\begin{tabular*}{\linewidth}{@{}@{\extracolsep{\fill}}llll@{}}
          \(\text{A. }a<b<c\) & \(\text{B. }b<a<c\) & \(\text{C. }c<b<a\) & \(\text{D. }c<a<b\)
        \end{tabular*}\par}

        \item $\left(x+\dfrac{1}{x}\right)\left(x+\dfrac{4}{x}\right)$ 最小值为
       
        {\centering\begin{tabular*}{\linewidth}{@{}@{\extracolsep{\fill}}llll@{}}
          \(\text{A. }10\) & \(\text{B. }9\) & \(\text{C. }8\) & \(\text{D. }6\)
        \end{tabular*}\par}

        \item 已知 $S_{2n}-S_n=n$，$a_3=6$，则 $a_7+a_8=$
        
        {\centering\begin{tabular*}{\linewidth}{@{}@{\extracolsep{\fill}}llll@{}}
          \(\text{A. }68\) & \(\text{B. }56\) & \(\text{C. }-3\) & \(\text{D. }-4\)
        \end{tabular*}\par}

        \item 已知双曲线 $\dfrac{x^2}{a^2}-\dfrac{y^2}{b^2}=1\,(a>0,b>0)$ 的左焦点为 $F$，$A$ 是右顶点．$P$ 是双曲线上一点，满足 $|FA|=|FP|$，$\angle FAP=30^\circ$，则双曲线离心率为
        
        {\centering\begin{tabular*}{\linewidth}{@{}@{\extracolsep{\fill}}llll@{}}
          \(\text{A. }4\) & \(\text{B. }\dfrac{8}{3}\) & \(\text{C. }\dfrac{8}{5}\) & \(\text{D. }\dfrac{4}{3}\)
        \end{tabular*}\par}
    \end{enumerate}
\end{description}

\begin{description}
    \item[二、填空题] 本题共6小题，每小题5分，共30分．
    \begin{enumerate}[leftmargin=5pt]
      \addtocounter{enumi}{+9}
        \item 化简 $(3+\mathrm{i})^2=$ \rule[-0.3ex]{3em}{0.4pt}．
        \item $(x+2y)^4$ 展开式中 $x^3y$ 的系数为 \rule[-0.3ex]{3em}{0.4pt}．
        \item 在 $\triangle ABC$ 中，$BC=4$，$AC=3$，$\cos A=-\dfrac{1}{4}$．则 $\sin B=$ \rule[-0.3ex]{3em}{0.4pt}．
        \item 箱子里有一个红球，两个黄球，三个白球．有放回地取三次，三次都没取到黄球的概率是 \rule[-0.3ex]{3em}{0.4pt}；在三次都没取到黄球的条件下，至少取到一次红球的概率是 \rule[-0.3ex]{3em}{0.4pt}．
        \item 已知 $|\vec{a}|=\vec{a}\cdot\vec{b}=1$，$|\vec{b}|>1$．记 $\vec{c}=\lambda\vec{a}+\mu\vec{b}$．当 $\vec{a}+\vec{b}-\vec{c}=\vec{0}$ 时，$\lambda+\mu=$ \rule[-0.3ex]{3em}{0.4pt}；当 $|\vec{a}+\vec{b}-\vec{c}|=1$ 时，$\lambda+\mu$ 范围 \rule[-0.3ex]{3em}{0.4pt}．
        \item 在平面内，抛物线 $y^2=2x$ 上有 $A,B,C,D$ 四个点．$A,B,C,D$ 的纵坐标分别为 $y_A,y_B,y_C,y_D$，直线 $AC$，直线 $BD$ 交 $x$ 轴于点 $P$，$AB$ 交 $x$ 轴于点 $M$，$CD$ 交 $x$ 轴于点 $N$．以下说法正确的有： \rule[-0.3ex]{3em}{0.4pt}．
        \begin{enumerate}[label=\textcircled{\arabic*}]
            \item 回忆版暂无
        \end{enumerate}
    \end{enumerate}
\end{description}

\begin{description}
    \item[三、解答题] 本题共5小题，共75分．解答应写出文字说明、证明过程或演算步骤．
    \begin{enumerate}[leftmargin=5pt]
      \addtocounter{enumi}{+15}
        \item （14分）

        $f(x)=\sin\left(2x+\dfrac{\pi}{6}\right)$．

        \begin{enumerate}[label=(\arabic*)]
            \item 求最小正周期；
            \item 若 $x\in\left[-\dfrac{\pi}{6},\dfrac{\pi}{12}\right]$，求 $f(x)$ 最大值和最小值；
            \item 若 $\alpha\in\left(0,\dfrac{\pi}{2}\right)$，$\sin\alpha=\dfrac{\sqrt{3}}{3}$，求 $f(\alpha)$．
        \end{enumerate}

        \item （15分）

        在长方体 $ABCD-A_1B_1C_1D_1$ 中，$AB=2$，$AA_1=3$，$AD=4$，$A_1E=3ED_1$，$C_1F=2FC$．

        \begin{enumerate}[label=(\arabic*)]
            \item 求证：$BD\perp$ 面 $CEF$．
            \item 求面 $AEF$ 与面 $CEF$ 的夹角的余弦值；
            \item 求三棱锥 $A-CEF$ 的体积．
        \end{enumerate}

        \item （15分）

        已知椭圆 $C:\dfrac{x^2}{a^2}+\dfrac{y^2}{b^2}=1\,(a>b>0)$ 的离心率为 $\dfrac{1}{2}$，椭圆被直线 $x=b$ 截得的线段长为 $\sqrt{3}$．

        \begin{enumerate}[label=(\arabic*)]
            \item 求 $C$ 的标准方程；
            \item 斜率为 $-\sqrt{3}$ 的直线与圆 $x^2+y^2=b^2$ 相切，且该直线交椭圆于 $P(x_1,y_1),Q(x_2,y_2)$，$y_1<y_2$．$A$ 是椭圆上顶点．记直线 $AP$，$AQ$ 的斜率分别为 $k_1,k_2$．求 $\dfrac{k_1}{k_2}$．
        \end{enumerate}

        \item （15分）

        已知等差数列 $\{a_n\}$ 与等比数列 $\{b_n\}$ 满足：$a_1=2$，$b_1=1$，$a_2=b_1+b_2$，$a_4=b_3-b_1$．

        \begin{enumerate}[label=(\arabic*)]
            \item 求数列 $\{a_n\}$，$\{b_n\}$ 的通项公式；
            \item 记 $E_n=\{x\in\mathbb{R}\mid x\leqslant n,\text{存在 }k\in\mathbb{N}^*\text{ 使得 }x=a_k\text{ 或 }x=b_k\}$，记 $c_n$ 为 $E_n$ 中的元素个数．
            \begin{enumerate}[label=(\roman*)]
                \item 求 $c_{3^m}$；
                \item 求 $\displaystyle\sum_{m=1}^{3^n-1}(-1)^m\cdot a_m\cdot c_m$．
            \end{enumerate}
        \end{enumerate}

        \item （16分）

        已知 $f(x)=e^x-\dfrac{2}{3}\sin x$．

        \begin{enumerate}[label=(\arabic*)]
            \item 求 $(0,f(0))$ 处的切线方程．
            \item 当 $x\in\left[-\dfrac{1}{3},+\infty\right)$ 时，证明 $f(x)\geqslant1+\dfrac{x}{3}$．
            \item 求实数 $\alpha$ 的最大可能值，使得 $f\!\left(\dfrac{1}{1}\right)\cdot f\!\left(\dfrac{1}{2}\right)\cdot f\!\left(\dfrac{1}{3}\right)\cdots f\!\left(\dfrac{1}{n}\right)\geqslant(n+1)^\alpha$ 对任意的 $n\in\mathbb{N}^*$ 都成立．
        \end{enumerate}
    \end{enumerate}
\end{description}
\clearpage

